\documentclass[11pt,a4paper]{article}

% --- Paketler ---
\usepackage[utf8]{inputenc}
\usepackage[T1]{fontenc}
% babel-turkish yoksa (texlive-lang-european) İngilizce'ye düşer
\IfFileExists{turkish.ldf}
    {\usepackage[shorthands=off,turkish]{babel}}
    {\usepackage[english]{babel}}
\usepackage{geometry}
\usepackage{xcolor}
\usepackage{titlesec}
\usepackage{enumitem}
\usepackage{booktabs}
\usepackage{array}
\usepackage{fancyhdr}
\usepackage{listings}
\usepackage{tcolorbox}
\usepackage{amsmath}
\usepackage{amssymb}
\usepackage{parskip}
\usepackage{hyperref}

% --- Sayfa Düzeni ---
\geometry{top=2.5cm, bottom=2.5cm, left=2.5cm, right=2.5cm}

% --- Renkler ---
\definecolor{primary}{RGB}{20, 60, 120}
\definecolor{secondary}{RGB}{200, 80, 20}
\definecolor{codebg}{RGB}{245, 245, 245}
\definecolor{notebg}{RGB}{255, 250, 230}
\definecolor{warnbg}{RGB}{255, 235, 235}
\definecolor{successbg}{RGB}{235, 255, 240}

% --- Başlık Formatları ---
\titleformat{\section}{\Large\bfseries\color{primary}}{\thesection}{1em}{}
\titleformat{\subsection}{\large\bfseries\color{secondary}}{\thesubsection}{1em}{}
\titleformat{\subsubsection}{\normalsize\bfseries}{\thesubsubsection}{1em}{}

% --- Kod Listeleme ---
% literate: pdflatex altında listings içindeki Türkçe karakterler için
\lstset{
    backgroundcolor=\color{codebg},
    basicstyle=\ttfamily\small,
    breaklines=true,
    frame=single,
    rulecolor=\color{gray},
    keywordstyle=\color{blue},
    commentstyle=\color{gray},
    stringstyle=\color{red},
    showstringspaces=false,
    tabsize=2,
    captionpos=b,
    inputencoding=utf8,
    extendedchars=true,
    literate={ç}{{\c{c}}}1 {Ç}{{\c{C}}}1 {ğ}{{\u{g}}}1 {Ğ}{{\u{G}}}1
             {ı}{{\i}}1 {İ}{{\.{I}}}1 {ö}{{\"o}}1 {Ö}{{\"O}}1
             {ş}{{\c{s}}}1 {Ş}{{\c{S}}}1 {ü}{{\"u}}1 {Ü}{{\"U}}1
             {→}{{$\rightarrow$}}1 {×}{{$\times$}}1
}

% --- tcolorbox Tanımları ---
\newtcolorbox{notebox}[1][]{colback=notebg, colframe=orange!70!black, title=#1,
    fonttitle=\bfseries, boxrule=0.5pt, arc=2mm}
\newtcolorbox{warnbox}[1][]{colback=warnbg, colframe=red!70!black, title=#1,
    fonttitle=\bfseries, boxrule=0.5pt, arc=2mm}
\newtcolorbox{successbox}[1][]{colback=successbg, colframe=green!50!black, title=#1,
    fonttitle=\bfseries, boxrule=0.5pt, arc=2mm}

% --- Header/Footer ---
\pagestyle{fancy}
\fancyhf{}
\fancyhead[L]{\small\textcolor{primary}{Fotonik AI Hızlandırıcı — Yol Haritası}}
\fancyhead[R]{\small\textcolor{primary}{\thepage}}
\renewcommand{\headrulewidth}{0.5pt}

% --- Hyperref Ayarları ---
\hypersetup{
    colorlinks=true,
    linkcolor=primary,
    urlcolor=secondary,
    citecolor=primary,
    pdftitle={Fotonik AI Hızlandırıcı Yol Haritası},
    pdfsubject={4 Yıllık Teknik ve Stratejik Plan}
}

% --- Başlık Bilgileri ---
\title{
    \vspace{-2cm}
    \Huge\textbf{\textcolor{primary}{Fotonik AI Hızlandırıcı}}\\[0.5cm]
    \Large\textbf{\textcolor{secondary}{4 Yıllık Uçtan Uca Yol Haritası}}\\[0.3cm]
    \large Dijital + Fotonik Hibrit Matris Çarpım Hızlandırıcısı\\[0.5cm]
    \normalsize\textit{Arty Z7-10 (Zynq-7010) Tabanlı Solo Geliştirme Planı}\\[0.2cm]
    \normalsize Revizyon 2 — Eylül 2026
}
\author{\textbf{Inference Engineer}}
\date{\today}

% ============================================================
\begin{document}

\maketitle
\thispagestyle{empty}

\vspace{1cm}

\begin{center}
\begin{tcolorbox}[colback=primary!10, colframe=primary, width=0.9\textwidth, arc=3mm]
\centering
\large\textit{``Ben bir çip fabrikası kurmuyorum.\\ Ben iki kişilik bir deep-tech stüdyosuyum.\\ Önce fikrin sayılarla işe yaradığını gösteririm,\\ sonra silikona para harcarım.''}
\end{tcolorbox}
\end{center}

\vspace{0.5cm}

\begin{notebox}[Revizyon 2'de Neler Değişti?]
\begin{itemize}[leftmargin=*,nosep]
    \item Donanım Arty Z7-20 değil, \textbf{Arty Z7-10 (xc7z010clg400-1)} olarak düzeltildi; PYNQ yerine ARM + AXI-Lite akışı.
    \item Sistem clock'u 125 MHz olarak düzeltildi (H16 pini).
    \item Efabless 2025'te kapandığı için ChipIgnite çıkarıldı, açık MPW alternatifleri eklendi.
    \item MZI matematiği düzeltildi (karmaşık 2$\times$2 unitary + SVD + Clements mesh).
    \item RISC-V çıkarıldı (Zynq'te zaten ARM var).
    \item Yıl 1'e \textbf{fotonik matmul gürültü/kuantizasyon simülatörü} ve \textbf{GPU baseline} eklendi.
    \item Her çeyreğe \textbf{go/no-go kriteri} eklendi.
    \item Kaynağı doğrulanamayan pazar/yetenek rakamları çıkarıldı.
    \item Gün gün plan, gerçekte tamamlanan işlerle (Hafta 1--2) değiştirildi.
\end{itemize}
\end{notebox}

\tableofcontents
\newpage

% ============================================================
\section{Vizyon ve Stratejik Konumlandırma}

\subsection{Neden Fotonik? Neden Şimdi?}

Modern yapay zeka sistemlerinde darboğaz giderek hesaplamadan \textbf{veri hareketine} kayıyor: bakır bağlantıların bant genişliği, gecikmesi ve bit başına enerjisi (bellek duvarı ve I/O darboğazı). Bu yüzden iki fotonik yön öne çıkıyor:

\begin{itemize}
    \item \textbf{Optik bağlantı (optical I/O, Co-Packaged Optics):} Çipler arası veri taşımayı ışıkla yapmak. Ticari olarak daha olgun.
    \item \textbf{Optik hesaplama:} Matris çarpımını girişim (interference) ile analog olarak yapmak. Daha riskli, daha az olgun.
\end{itemize}

\begin{notebox}[Pazar Verisi — Kaynak Eklenecek]
Fon başvurularında kullanılacak pazar büyüklüğü rakamları (CPO pazar tahminleri vb.) \textbf{birincil kaynaktan} (analist raporu, şirket duyurusu) alınıp tarih ve bağlantısıyla buraya eklenmeli. Doğrulanamayan bir rakam, başvurunun geri kalanının güvenilirliğini de düşürür.
\end{notebox}

\subsection{Rekabet Ortamı}

\begin{warnbox}[Uyarı: Rekabet Yoğun]
Bu alan ``boş arazi'' değil, iyi fonlanmış şirketlerin olduğu bir arena. Kazanma şansın \textbf{odaklı, çevik ve niş} olmaktan geçiyor.
\end{warnbox}

\begin{table}[h!]
\centering
\small
\begin{tabular}{@{}p{3cm}p{2cm}p{9cm}@{}}
\toprule
\textbf{Şirket} & \textbf{Bölge} & \textbf{Odak} \\
\midrule
Lightmatter & ABD & Fotonik bağlantı (interconnect) ve fotonik hesaplama \\
Ayar Labs & ABD & Optik I/O chiplet'leri \\
Lightelligence & Çin & Optik hesaplama ve optik bağlantı \\
Neurophos & ABD & Metasurface tabanlı optik hesaplama \\
\bottomrule
\end{tabular}
\caption{Fotonik AI alanındaki bazı oyuncular. Fon ve değerleme rakamları birincil kaynaktan doğrulanıp eklenecek.}
\end{table}

\subsection{Senin Nişin}

\textbf{Onların yaklaşımı:} Genel amaçlı, büyük ölçekli optik işlemciler veya bağlantı çözümleri.

\textbf{Senin fırsatın:} Inference mühendisliği geçmişin (GPU inference, GenAI görüntü/video modelleri) sana şu soruyu sorduruyor:

\begin{center}
\large\textit{``Tüm modeli optik yapmak zorunda mıyız? Yoksa sadece en çok enerji tüketen belirli bir katmanı (örneğin Attention veya belirli bir MatMul) optik olarak hızlandırsak yeterli olmaz mı?''}
\end{center}

Bu soru \textbf{önce simülasyonla, sayılarla} cevaplanmalı (Bölüm~\ref{sec:simulator}). Cevap olumluysa donanıma yatırım yapılır; olumsuzsa hedef erkenden, ucuza değiştirilir.

% ============================================================
\section{Durum Analizi ve Varlıklar}

\subsection{Elimizdeki Donanım: Arty Z7-10}

\begin{table}[h!]
\centering
\small
\begin{tabular}{@{}p{5.2cm}p{9cm}@{}}
\toprule
\textbf{Özellik} & \textbf{Senin İçin Anlamı} \\
\midrule
Zynq-7010 SoC (çift çekirdek ARM Cortex-A9 + FPGA) & Yazılım (C) ile donanımı (RTL) aynı çipte, AXI üzerinden birleştirebilirsin. \\
FPGA: 17.600 LUT, 35.200 FF & Küçük kontrol mantığı ve wrapper'lar için yeterli. \\
80 DSP48E1 & 8$\times$8 INT8 systolic array'e kadar (64 DSP) sığar; 4$\times$4 dizi 16 DSP kullanıyor. \\
60 BRAM36 ($\sim$2,1 Mb) & Matris tile'ları için yerel tampon. \\
512 MB DDR3 (16-bit) & ARM tarafı için ana bellek. \\
125 MHz PL clock (H16 pini) & XDC'de \texttt{-period 8.00} olmalı. \\
XADC (12-bit, 1 MSPS) & Fotodetektörden yavaş/DC ölçüm ve kalibrasyon için yeterli; gerçek zamanlı yüksek hızlı okuma için harici ADC gerekir. \\
\bottomrule
\end{tabular}
\caption{Arty Z7-10'un projeye katkısı}
\end{table}

\begin{warnbox}[PYNQ Z7-10'da Resmi Olarak Desteklenmiyor]
Resmi PYNQ imajları PYNQ-Z1/Z2 gibi Zynq-7020 kartları içindir; Arty Z7-10'daki 7010 için hazır imaj yok. Bu yüzden plan \textbf{Vivado block design (ZYNQ7 PS + AXI-Lite) + Vitis bare-metal C} akışına göre yazıldı. Linux + Python gerekirse ileride PetaLinux veya topluluk PYNQ port'u değerlendirilebilir.
\end{warnbox}

\begin{notebox}[DAC Eksikliği]
Kartta XADC (giriş) var ama \textbf{DAC (çıkış) yok}. Fotonik çipteki faz kaydırıcıları sürmek için harici bir DAC gerekecek. Termal faz kaydırıcılar yavaş (DC'ye yakın) sürüldüğü için ilk aşamada çok kanallı, düşük hızlı bir DAC (ör. Pmod DA4) yeterli; yüksek hızlı modülatörler için ayrı bir çözüm gerekir.
\end{notebox}

\subsection{Takım Yapısı}

\begin{table}[h!]
\centering
\small
\begin{tabular}{@{}p{2cm}p{6cm}p{6cm}@{}}
\toprule
\textbf{Dönem} & \textbf{Sen (Solo)} & \textbf{Kuzenin (Yarı Zamanlı)} \\
\midrule
1. yıl & RTL + systolic array + AXI + simülatör + Tiny Tapeout & Dersler + yaz tatilinde gdsfactory/SAX \\
2. yıl & DAC/ADC arayüzü + ASIC (açık MPW) & Fotonik dersleri + MZI mesh simülasyonu \\
3. yıl & Sürücü ASIC + hibrit mimari & İlk PIC layout + üniversite laboratuvarı \\
4. yıl & Sistem entegrasyonu + demo + fon & Mezuniyet projesi + tam zamanlı katılım \\
\bottomrule
\end{tabular}
\caption{4 yıllık rol dağılımı}
\end{table}

% ============================================================
\section{4 Yıllık Yol Haritası}

\subsection{Yıl 1 (Eylül 2026 -- Ağustos 2027): Dijital Temel + Fikrin Doğrulanması}

\begin{table}[h!]
\centering
\small
\begin{tabular}{@{}p{1.6cm}p{5.2cm}p{5cm}p{2.4cm}@{}}
\toprule
\textbf{Çeyrek} & \textbf{Dijital} & \textbf{Fotonik / Simülasyon} & \textbf{Çıktı} \\
\midrule
Q1 \newline Eyl--Kas & 4$\times$4 systolic array + cocotb (\checkmark), AXI-Lite wrapper, ARM'dan test & gdsfactory + SAX ile ilk MZI, 2$\times$2 transfer matrisi & Kartta çalışan matmul \\
Q2 \newline Ara--Şub & Seri yüklemeli dizi (Tiny Tapeout I/O'suna uygun), GPU baseline (CuTe/CUTLASS) & Gürültü + kuantizasyon simülatörü, Clements mesh ile SVD matmul & Simülatör raporu \\
Q3 \newline Mar--May & XADC + DAC arayüz prototipi & Gerçek bir katmanı (attention/MLP) simülatörde çalıştır & \textbf{Go/no-go kararı} \\
Q4 \newline Haz--Ağu & \textbf{Tiny Tapeout başvurusu} & PIC mimari spesifikasyonu & İlk dijital çip tasarımı \\
\bottomrule
\end{tabular}
\caption{Yıl 1 yol haritası}
\end{table}

\subsection{Yıl 2 (2027--2028): Fotonik Tarafa Geçiş}

\begin{itemize}
    \item Tiny Tapeout çipi geldiğinde test et, test PCB'si tasarla.
    \item DAC/ADC + fotonik arayüz mimarisini kur.
    \item \textbf{İlk ASIC tape-out} (açık MPW; bkz. Bölüm~\ref{sec:mpw}).
    \item \textbf{İlk PIC tape-out} (fotonik MPW; bkz. Bölüm~\ref{sec:mpw}).
\end{itemize}

\subsection{Yıl 3 (2028--2029): Entegrasyon}

\begin{itemize}
    \item Sürücü çipini test et, hataları düzelt, ikinci revizyonu planla.
    \item PIC'i test et, optik kayıpları ölç, kalibrasyon rutinleri yaz.
    \item İki çipi aynı PCB'de birleştirme denemesi (MCM).
\end{itemize}

\subsection{Yıl 4 (2029--2030): Demo + Fon}

\begin{itemize}
    \item Küçük bir DiT/UNet katmanını hibrit sistemde çalıştır.
    \item Yıl 1'de ölçülen GPU baseline'ına karşı benchmark yap.
    \item TÜBİTAK BİGG, EIC Accelerator, YC gibi fonlara başvur.
    \item Kuzenin mezun olur, tam zamanlı katılır.
\end{itemize}

\subsection{Tape-Out Seçenekleri}
\label{sec:mpw}

\begin{warnbox}[Efabless Kapandı]
Efabless 2025'te faaliyetlerini durdurdu; \textbf{ChipIgnite artık bir seçenek değil}. Aşağıdaki seçeneklerin güncel durumu, fiyatı ve takvimi başvurudan önce kontrol edilmeli.
\end{warnbox}

\begin{table}[h!]
\centering
\small
\begin{tabular}{@{}p{3.5cm}p{3cm}p{7.5cm}@{}}
\toprule
\textbf{Seçenek} & \textbf{Tür} & \textbf{Not} \\
\midrule
Tiny Tapeout & Dijital, çok küçük & Yıl 1 hedefi. Az I/O pini, küçük alan, sınırlı I/O hızı. \\
IHP açık MPW & Dijital (SG13G2) & Açık PDK, akademik/açık kaynak projeler için. \\
wafer.space & Dijital (GF180MCU) & Düşük maliyetli açık PDK çalıştırmaları. \\
Europractice & Dijital + fotonik & Akademik erişim; üniversite üzerinden daha uygun. \\
Cornerstone (UK) & Silikon fotonik & Akademik odaklı fotonik MPW. \\
AIM Photonics (ABD) & Silikon fotonik & Daha pahalı, olgun PDK. \\
\bottomrule
\end{tabular}
\caption{Tape-out seçenekleri (güncel durumu kontrol edilecek)}
\end{table}

% ============================================================
\section{Go/No-Go Kriterleri}

Solo çalışmanın en büyük riski, yanlış bir yolda aylarca ilerlemek. Her çeyrek sonunda aşağıdaki sorulara \textbf{ölçülmüş sayılarla} cevap verilir. Eşik değerleri başlangıç önerisidir; simülatör sonuçlarına göre güncellenir.

\begin{table}[h!]
\centering
\small
\begin{tabular}{@{}p{1.6cm}p{6.4cm}p{6cm}@{}}
\toprule
\textbf{Ne zaman} & \textbf{Kriter} & \textbf{Karşılanmazsa} \\
\midrule
Q1 sonu & ARM'dan gönderilen rastgele 4$\times$4 INT8 matrisler, yazılım referansıyla \%100 aynı sonucu veriyor & Q2'ye geçme, hata ayıkla \\
Q2 sonu & Simülatör, bilinen sonuçları (ideal durumda hatasız matmul) tekrar üretiyor & Model doğrulanmadan sonuçlara güvenme \\
Q3 sonu & Hedef katmanda, gerçekçi analog hassasiyetle (ör. 6 bit) doğruluk kaybı kabul edilebilir (ör. \textless\%1) \textbf{ve} tahmini pJ/MAC, GPU baseline'ından daha iyi & Hedef katmanı veya optik hedefini değiştir (ör. optik bağlantıya odaklan) \\
Q4 sonu & Tiny Tapeout tasarımı gate-level simülasyonu ve tüm kontrolleri geçiyor & Bir sonraki shuttle'a kaydır \\
\bottomrule
\end{tabular}
\caption{Çeyreklik go/no-go kriterleri}
\end{table}

% ============================================================
\section{Tamamlananlar ve Sıradaki Adımlar}

\subsection{Hafta 1 (Eylül 2026): Ortam + İlk RTL}

\begin{itemize}
    \item Vivado 2026.1 (Linux) kuruldu, Digilent board dosyaları eklendi.
    \item Kartın Z7-20 değil \textbf{Z7-10} olduğu JTAG IDCODE ile anlaşıldı; proje yeniden oluşturuldu.
    \item \texttt{led\_blink.sv} sentez $\rightarrow$ implementation $\rightarrow$ bitstream $\rightarrow$ FPGA: LED yanıp sönüyor.
    \item FTDI udev kuralı ve GitHub token sorunları çözüldü.
\end{itemize}

\subsection{Hafta 2: 4$\times$4 INT8 Systolic Array}

\begin{itemize}
    \item PE modülü + cocotb testi (Verilator).
    \item 4$\times$4 output-stationary systolic array; 16/16 sonuç doğru.
    \item \texttt{use\_dsp} attribute'u sinyal bildirimine taşındı: her PE tek bir DSP48E1 MAC'e oturdu (16 DSP, 0 LUT, 8 FF).
    \item Out-of-context sentez: IOB sorunu olmadan kaynak ve timing tahmini; 100 MHz'de WNS +7,995 ns (I/O yolları henüz kısıtlanmadı).
\end{itemize}

\textbf{Güncel PE modülü} (hedef: akümülatör 32 bit):

\begin{lstlisting}[language=Verilog]
module pe #(
    parameter int WIDTH     = 8,
    parameter int ACC_WIDTH = 32   // 16 bit K>1 için taşabilir
) (
    input  logic                        clk,
    input  logic                        rst,
    input  logic                        en,
    input  logic signed [WIDTH-1:0]     a_in,
    input  logic signed [WIDTH-1:0]     b_in,
    output logic signed [WIDTH-1:0]     a_out,
    output logic signed [WIDTH-1:0]     b_out,
    output logic signed [ACC_WIDTH-1:0] acc
);
    (* use_dsp = "yes" *) logic signed [2*WIDTH-1:0] mult;
    assign mult = a_in * b_in;

    always_ff @(posedge clk) begin
        if (rst) begin
            a_out <= '0;
            b_out <= '0;
            acc   <= '0;
        end else if (en) begin
            a_out <= a_in;
            b_out <= b_in;
            acc   <= acc + mult;
        end
    end
endmodule
\end{lstlisting}

\begin{warnbox}[Akümülatör Taşması]
INT8$\times$INT8 çarpımın en büyük değeri $(-128)\times(-128)=16384$. K=4 için toplam 65536'ya ulaşabilir; bu int16 sınırını (32767) aşar. Gerçek INT8 hızlandırıcılar INT32 akümülatör kullanır. DSP48E1'in P çıkışı zaten 48 bit olduğu için ek kaynak gerekmez.
\end{warnbox}

\subsection{Hafta 3--4: Sıradaki Adımlar}

\begin{enumerate}
    \item \textbf{Küçük düzeltmeler:} XDC'de clock'u 125 MHz yap, \texttt{acc}'yi 32 bit yap, cocotb testlerine en kötü durum (taşma) testi ekle.
    \item \textbf{I/O kısıtları:} OOC akışında \texttt{set\_input\_delay}/\texttt{set\_output\_delay} ekleyip implementation sonrası gerçek timing'i al.
    \item \textbf{AXI-Lite wrapper:} A (16 byte), B (16 byte), start/done ve C (16$\times$4 byte) register haritası.
    \item \textbf{Block design:} ZYNQ7 Processing System + AXI-Lite wrapper; bitstream + XSA export.
    \item \textbf{Vitis bare-metal test:} ARM'dan rastgele matrisler yaz, C'yi oku, yazılım referansıyla karşılaştır.
    \item \textbf{Fotonik hat:} \texttt{photonic/} klasörünü aç, gdsfactory ile ilk MZI, SAX ile transfer matrisi.
\end{enumerate}

\textbf{LED blink (125 MHz clock için düzeltilmiş):}

\begin{lstlisting}[language=Verilog]
module led_blink (
    input  logic clk,   // 125 MHz, H16
    input  logic rst,
    output logic led
);
    // 125e6 / 2 = 62_500_000 -> 0,5 s'de bir toggle (1 Hz yanıp sönme)
    localparam int unsigned HALF_PERIOD = 62_500_000;
    logic [26:0] counter;

    always_ff @(posedge clk) begin
        if (rst) begin
            counter <= '0;
            led     <= 1'b0;
        end else if (counter == HALF_PERIOD - 1) begin
            counter <= '0;
            led     <= ~led;
        end else begin
            counter <= counter + 1;
        end
    end
endmodule
\end{lstlisting}

% ============================================================
\section{Fotonik Matmul Simülatörü}
\label{sec:simulator}

Yıl 1'in en önemli çıktısı. Amaç, tape-out'a para harcamadan önce şu soruları \textbf{sayılarla} cevaplamak:

\begin{itemize}
    \item Gerçekçi analog hassasiyetle (faz hatası, dedektör gürültüsü, DAC/ADC bit derinliği) hedef katmanın doğruluğu ne kadar düşüyor?
    \item DAC/ADC ve lazer dahil toplam enerji, GPU'daki aynı GEMM'e göre pJ/MAC olarak nerede?
    \item Hangi katman optik, hangisi dijital kalmalı?
\end{itemize}

\begin{notebox}[Neden Önemli?]
Fotonik hesaplamada zorluk çoğu zaman optik çarpmanın kendisinde değil, \textbf{elektro-optik dönüşümde}: DAC/ADC enerjisi ve sınırlı analog hassasiyet. Bu yüzden simülatör DAC/ADC'yi de modellemeli.
\end{notebox}

\textbf{İlk yaklaşım (PyTorch):}

\begin{lstlisting}[language=Python]
import torch

def quantize(x, bits, xmax):
    # Simetrik uniform kuantizasyon (DAC/ADC modeli)
    levels = 2 ** (bits - 1) - 1
    step = xmax / levels
    return torch.clamp(torch.round(x / step), -levels, levels) * step

def photonic_matmul(x, W, dac_bits=6, adc_bits=6,
                    weight_noise=0.01, det_noise=0.01):
    # 1) Giriş DAC ile kodlanır
    x_q = quantize(x, dac_bits, x.abs().max())
    # 2) Ağırlıklar faz hatası nedeniyle bozulur (birinci derece yaklaşım;
    #    sonraki adım: Clements mesh fazlarına doğrudan gürültü eklemek)
    W_n = W * (1 + weight_noise * torch.randn_like(W))
    y = x_q @ W_n.T
    # 3) Fotodetektör gürültüsü
    y = y + det_noise * y.abs().max() * torch.randn_like(y)
    # 4) Çıkış ADC ile okunur
    return quantize(y, adc_bits, y.abs().max())
\end{lstlisting}

Sonraki adımlar: ağırlık gürültüsünü doğrudan MZI fazlarında modellemek (Clements mesh), optik kayıp ve lazer gücünü eklemek, bir DiT attention/MLP katmanında \texttt{torch.nn.Linear}'ı bu fonksiyonla değiştirip doğruluk kaybını ölçmek.

\subsection{GPU Baseline}

Aynı GEMM boyutları için CuTe/CUTLASS ile GPU'da latency, throughput ve (mümkünse güç ölçümüyle) pJ/MAC değerleri Yıl 1 Q2'de ölçülür. Bu, Yıl 4 benchmark'ının referans çizgisidir.

% ============================================================
\section{Fotonik Öğrenme Yol Haritası}

\subsection{Inference Engineer İçin Eşleştirme}

\begin{table}[h!]
\centering
\small
\begin{tabular}{@{}p{5cm}p{9cm}@{}}
\toprule
\textbf{Bildiğin (AI/Inference)} & \textbf{Fotonikte Karşılığı} \\
\midrule
MatMul (GEMM) & MZI mesh ile analog matris çarpımı (SVD) \\
Quantization (INT8/FP8) & DAC/ADC bit derinliği, faz kaydırıcı çözünürlüğü \\
Bellek duvarı & Elektro-optik dönüşüm maliyeti, optik kayıp \\
Kernel fusion & Fotonik bileşenleri tek çipte birleştirme (PIC) \\
Graph partitioning & Hangi katman optik, hangisi dijital? \\
Latency/throughput & Optik gecikme vs elektronik gecikme \\
Enerji bütçesi & Lazer gücü, modülatör enerjisi, fotodetektör gürültüsü \\
\bottomrule
\end{tabular}
\caption{AI bilgisinden fotonik bilgisine köprü}
\end{table}

\subsection{MZI Matmul Matematiği}

Bir MZI, iki faz kaydırıcısı ($\theta$: iç faz, $\phi$: dış faz) ile iki dalga kılavuzu arasında \textbf{karmaşık} bir 2$\times$2 unitary dönüşüm uygular. Clements et al. (2016) gösterimiyle:

\begin{equation}
T(\theta, \phi) = \begin{bmatrix}
e^{i\phi}\cos\theta & -\sin\theta \\
e^{i\phi}\sin\theta & \cos\theta
\end{bmatrix}
\end{equation}

$N(N-1)/2$ adet MZI'dan oluşan bir mesh (Reck veya Clements dizilimi), herhangi bir $N\times N$ unitary matrisi $U$'yu gerçekleyebilir. Ancak AI ağırlık matrisleri genelde unitary değildir. Keyfi bir $W$ matrisi için \textbf{tekil değer ayrışımı (SVD)} kullanılır:

\begin{equation}
W = U \Sigma V^{\dagger}
\end{equation}

\begin{itemize}
    \item $V^{\dagger}$ ve $U$: iki ayrı MZI mesh ile gerçeklenir.
    \item $\Sigma$: köşegen matris; her kanalda bir zayıflatıcı (veya kazanç) ile gerçeklenir.
\end{itemize}

Bu, Shen et al. (2017) çalışmasının temelidir. Pratik sorunlar: faz hatası mesh boyunca birikir, optik kayıp $N$ ile artar ve dedektör yalnızca yoğunluk ($|E|^2$) ölçer; işaretli değerler için ek kodlama gerekir.

\subsection{İlk MZI (gdsfactory)}

\begin{lstlisting}[language=Python]
import gdsfactory as gf

c = gf.components.mzi(delta_length=10)
c.write_gds("mzi_demo.gds")
c.plot()
\end{lstlisting}

\subsection{Araç Zinciri}

\begin{table}[h!]
\centering
\small
\begin{tabular}{@{}p{3.2cm}p{6.3cm}p{4.5cm}@{}}
\toprule
\textbf{Araç} & \textbf{Ne İçin?} & \textbf{Kurulum} \\
\midrule
gdsfactory & Python tabanlı fotonik layout & \texttt{pip install gdsfactory} \\
SAX & Fotonik devre (S-parametre) simülasyonu & \texttt{pip install sax} \\
Meep & FDTD simülasyonu & \texttt{conda install -c conda-forge pymeep} \\
KLayout & GDS görüntüleme ve DRC & klayout.de \\
PyTorch & Gürültü/kuantizasyon simülatörü & \texttt{pip install torch} \\
Lumerical & Ticari FDTD/devre simülasyonu & Üniversite lisansı \\
\bottomrule
\end{tabular}
\caption{Fotonik araç zinciri}
\end{table}

\subsection{Öğrenme Kaynakları}

\begin{itemize}
    \item L. Chrostowski, M. Hochberg, \textit{Silicon Photonics Design}, Cambridge University Press, 2015 (ve UBCx/edX kursu).
    \item Y. Shen et al., ``Deep learning with coherent nanophotonic circuits'', \textit{Nature Photonics}, 2017.
    \item W. R. Clements et al., ``Optimal design for universal multiport interferometers'', \textit{Optica}, 2016.
    \item M. Reck et al., ``Experimental realization of any discrete unitary operator'', \textit{Physical Review Letters}, 1994.
    \item gdsfactory: \url{https://gdsfactory.github.io/gdsfactory/}
    \item SAX: \url{https://flaport.github.io/sax/}
    \item Meep: \url{https://meep.readthedocs.io}
    \item SiEPIC (açık kaynak fotonik tasarım araçları ve eğitimleri)
    \item Tiny Tapeout: \url{https://tinytapeout.com}
\end{itemize}

% ============================================================
\section{Kariyer Güvenliği}

Silikon fotonik, AI veri merkezlerinin itmesiyle büyüyen bir alan ve \textbf{dijital tasarım + fotonik + AI inference} bileşimine sahip mühendis sayısı az. Bu proje, girişim başarılı olmasa bile bu bileşimi somut çıktılarla (RTL, tape-out, simülatör, benchmark) belgelemeni sağlar.

\begin{table}[h!]
\centering
\small
\begin{tabular}{@{}p{3.5cm}p{5cm}p{5cm}@{}}
\toprule
\textbf{Senaryo} & \textbf{Ne Olur?} & \textbf{Kariyer Sonucu} \\
\midrule
Girişim başarılı & Çip çalışır, fon bulunur & Kurucu/CTO \\
Kısmen başarılı & Demo çalışır, ticarileşemez & Acquihire, senior mühendis \\
Başarısız & Teknik hedefler tutmaz & Güçlü portfolyo ile büyük şirkette senior rol \\
\bottomrule
\end{tabular}
\caption{Üç senaryo}
\end{table}

\begin{notebox}[Doğrulama Notu]
Yetenek açığı ve maaş verileri fon başvurusunda kullanılacaksa, güncel ve kaynaklı veriyle (iş ilanları, sektör raporları) eklenmeli.
\end{notebox}

% ============================================================
\section{Kritik Uyarılar ve Prensipler}

\begin{warnbox}[IP Duvarı]
Mevcut işin sana GPU mimarisi ve inference deneyimi ile maaş runway'i veriyor. Ama işverenin kodunu, modelini, altyapısını kullanma; iş bilgisayarında bu proje üzerinde çalışma. \textbf{Clean-room çalış.} Dışarıyla paylaşılan dokümanlarda işveren adı ve iç bilgiler yer almasın.
\end{warnbox}

\begin{warnbox}[İşini Hemen Bırakma]
Akşamları ve hafta sonları bu projeye ayır. 12--18 aylık finansal plan yapmadan tam zamanlı girişimcilik seni yakar.
\end{warnbox}

\begin{warnbox}[Ortaklık Sözleşmesi]
Kuzeninle ortaklık sözleşmesini şimdiden konuş: hisse, vesting, IP, ayrılık durumu. Aile içi ortaklıkta en büyük risk duygusal varsayımlardır.
\end{warnbox}

\begin{notebox}[Tiny Tapeout Kısıtları]
Tiny Tapeout'ta proje başına az sayıda I/O pini (8 giriş, 8 çıkış, 8 çift yönlü), küçük alan ve sınırlı I/O hızı var. 4$\times$4 dizi FPGA'da 323 I/O istiyordu; aynı sorun orada da çıkar. Tasarım baştan \textbf{seri yükleme} (A/B'yi byte byte yükle, C'yi byte byte oku) ile yapılmalı.
\end{notebox}

\begin{notebox}[Kapsam Kayması]
``Bir de şunu ekleyeyim'' tuzağına düşme. Her çeyrekte \textbf{tek bir somut çıktı} ve bir go/no-go kriteri belirle.
\end{notebox}

\begin{notebox}[Tıkanma Noktaları]
1 haftadan fazla takılma. Discord (Tiny Tapeout, gdsfactory), GitHub Issues, hocalar.
\end{notebox}

% ============================================================
\section{Bu Hafta Yapacağın 3 Şey}

\begin{enumerate}
    \item \textbf{Düzeltmeler:} XDC clock'u 125 MHz, \texttt{acc} 32 bit, taşma testi.
    \item \textbf{AXI-Lite wrapper:} Register haritasını yaz, cocotb ile test et.
    \item \textbf{Fotonik hattı başlat:} \texttt{photonic/} klasörü, ilk MZI, SAX transfer matrisi.
\end{enumerate}

\begin{center}
\begin{tcolorbox}[colback=successbg, colframe=green!50!black, width=0.9\textwidth, arc=3mm]
\centering
\large\textbf{4 yıl sonra hedef:}\\[0.3cm]
\textbf{FPGA'da doğrulanmış, silikona dönüşmüş ve ışıkla çarpan bir çip.}\\[0.3cm]
\textit{İlk \texttt{always\_ff} bloğu yazıldı, 16 DSP'ye oturdu. Sıradaki: ARM'dan konuşmak.}
\end{tcolorbox}
\end{center}

% ============================================================
\appendix

\section{Ek: Hızlı Referans Komutları}

\subsection{Vivado Proje Akışı (Arty Z7-10)}

\begin{lstlisting}[language=bash]
# 1. Create Project -> RTL Project
# 2. Boards -> Arty Z7-10 (xc7z010clg400-1)
# 3. Add Sources -> rtl/*.sv, Add Constraints -> *.xdc
# 4. Run Synthesis -> Run Implementation -> Generate Bitstream
# 5. Open Hardware Manager -> Auto Connect -> Program Device
\end{lstlisting}

\subsection{Out-of-Context Sentez ve Timing (Tcl Console)}

\begin{lstlisting}[language=tcl]
set_property -name {STEPS.SYNTH_DESIGN.ARGS.MORE OPTIONS} \
    -value {-mode out_of_context} -objects [get_runs synth_1]
reset_run synth_1
launch_runs impl_1 -jobs 4
wait_on_run impl_1
open_run impl_1
report_utilization
report_timing_summary -max_paths 5
\end{lstlisting}

\subsection{Arty Z7 Clock Kısıtı}

\begin{lstlisting}[language=tcl]
set_property -dict { PACKAGE_PIN H16 IOSTANDARD LVCMOS33 } [get_ports { clk }]
create_clock -add -name sys_clk_pin -period 8.00 -waveform {0 4} [get_ports { clk }]
\end{lstlisting}

\subsection{cocotb Testleri}

\begin{lstlisting}[language=bash]
cd fpga/systolic_array_4x4/tb
make                                   # 4x4 dizi testi
make TOPLEVEL=pe MODULE=test_pe        # tek PE testi
\end{lstlisting}

\subsection{GitHub Repo Yapısı}

\begin{lstlisting}
photonic-inference-lab/
|-- fpga/
|   |-- led_blink_z010/          # Vivado projesi (Hafta 1)
|   |-- systolic_array_4x4/
|   |   |-- rtl/                 # pe.sv, systolic_array_4x4.sv
|   |   `-- tb/                  # cocotb testleri
|   `-- systolic_4x4_synth/      # Vivado sentez projesi
|-- photonic/                    # (eklenecek) gdsfactory, SAX, Meep
|-- sim/                         # (eklenecek) fotonik matmul simülatörü
|-- docs/
|   |-- roadmap.tex
|   `-- week_logs/
`-- README.md
\end{lstlisting}

\section{Ek: 4 Yıllık Özet Takvim}

\begin{table}[h!]
\centering
\small
\begin{tabular}{@{}p{0.8cm}p{3.4cm}p{3.4cm}p{3.4cm}p{2.6cm}@{}}
\toprule
\textbf{Yıl} & \textbf{Dijital Hat} & \textbf{Fotonik Hat} & \textbf{Ortak Çıktı} & \textbf{Fon} \\
\midrule
1 & Systolic array + AXI + Tiny Tapeout & MZI + gürültü simülatörü + GPU baseline & Go/no-go kararı + ilk dijital çip & TÜBİTAK 2209 / 1512 \\
2 & DAC/ADC + sürücü ASIC (açık MPW) & İlk PIC tape-out + kalibrasyon & İki çip + test PCB & TÜBİTAK BİGG \\
3 & Sürücü ASIC test + revizyon & Optik ölçüm + paketleme & Hibrit modül & KOSGEB Ar-Ge \\
4 & Sistem entegrasyonu + benchmark & Demo & Çalışan hibrit matmul demosu & EIC Accelerator \\
\bottomrule
\end{tabular}
\caption{4 yıllık özet takvim (fon programlarının uygunluk koşulları kontrol edilecek)}
\end{table}

\section{Ek: Sözlük}

\begin{description}
    \item[PIC] Photonic Integrated Circuit — Fotonik Entegre Devre
    \item[MZI] Mach-Zehnder Interferometer — Mach-Zehnder Girişimölçeri
    \item[SVD] Singular Value Decomposition — Tekil Değer Ayrışımı
    \item[CPO] Co-Packaged Optics — Ko-Paketlenmiş Optik
    \item[MPW] Multi-Project Wafer — Çoklu Proje Gofreti
    \item[PDK] Process Design Kit — Süreç Tasarım Kiti
    \item[OOC] Out-of-Context — I/O buffer'sız, modül seviyesinde sentez
    \item[RTL] Register Transfer Level — Kayıt Transfer Seviyesi
    \item[GDSII] Graphic Data System II — Layout dosya formatı
    \item[STA] Static Timing Analysis — Statik Zaman Analizi
    \item[WNS] Worst Negative Slack — En kötü zamanlama marjı
    \item[DRC] Design Rule Check — Tasarım Kuralı Kontrolü
    \item[LVS] Layout Versus Schematic — Layout-Şema Uyumu
    \item[XADC] Xilinx Analog-to-Digital Converter
    \item[AXI] Advanced eXtensible Interface — ARM ile FPGA arasındaki veri yolu
    \item[FDTD] Finite-Difference Time-Domain — Sonlu Farklar Zaman Alanı
    \item[GEMM] General Matrix Multiply — Genel Matris Çarpımı
\end{description}

\end{document}
